% =====================================================================
%   §5. Хеширование, теория чисел и рандомизация
% =====================================================================
\section{Хеширование: теория чисел встречается с рандомизацией}
\label{sec:hashing}

В предыдущих параграфах мы видели, как теория чисел даёт «жёсткие» гарантии безопасности: вычислить нужный ответ практически невозможно из-за вычислительной сложности задачи факторизации. Сейчас мы рассмотрим совсем другой тип использования теории чисел --- \term{вероятностные алгоритмы}, в которых правильность работы достигается не «навсегда», а \emph{в среднем}, с очень высокой вероятностью.

Главная мысль: добавив в алгоритм \emph{случайный выбор} (на этапе настройки, ещё до запуска), можно гарантировать, что для \emph{любых} входных данных алгоритм работает в среднем хорошо. Принципиальное отличие от наивного «алгоритм работает хорошо для случайных данных» в том, что входные данные могут быть какими угодно --- даже подобраны злоумышленником, --- а случайность спрятана \emph{внутри} алгоритма.

Этот параграф следует изложению из классической книги С.\ Дасгупты, Х.\ Пападимитриу и У.\ Вазирани «Алгоритмы»\footnote{Дасгупта~С., Пападимитриу~Х., Вазирани~У. \emph{Алгоритмы}. М.: МЦНМО, 2014. Глава о хешировании.}, в котором роль теории чисел особенно прозрачна.

% ---------------------------------------------------------------------
\subsection{Задача быстрого поиска}

Представьте сервис электронной почты, в котором зарегистрировано сто миллионов пользователей. Каждый пользователь идентифицируется адресом, скажем, длиной до 30 символов. При каждом входе нужно \emph{мгновенно} (за время, не зависящее от размера базы) определить: есть ли такой пользователь, и если есть --- найти его данные.

Технически: имеется большая «вселенная» возможных ключей $U$ (всех допустимых адресов), и интересующее нас подмножество $S \subseteq U$ из $n$ реальных пользователей. Мы хотим организовать структуру данных, позволяющую за $O(1)$ операций проверять, лежит ли $x \in S$, и получать данные, связанные с $x$.

\paragraph*{Прямая адресация.} Самое простое решение --- завести массив, индексируемый \emph{всеми} элементами $U$. Такой массив имеет длину $|U|$. Если $|U|$ велико (а оно \emph{колоссально}: количество всевозможных 30-символьных строк --- это $26^{30}$, что больше $10^{42}$), это решение нереализуемо: не хватит памяти всей планеты.

\paragraph*{Идея хеширования.} Вместо того, чтобы выделять ячейку под каждый возможный элемент, заведём массив размера $m \approx n$ (того же порядка, что число реальных пользователей), и зададим \term{хеш-функцию} $h: U \to \{0, 1, \dots, m-1\}$, которая каждому ключу сопоставляет номер ячейки.

\begin{definition}{хеш-функция и хеш-таблица}{def:p5-1}
\term{Хеш-функция} --- это произвольная функция $h: U \to \{0, 1, \dots, m-1\}$. По ней строится \term{хеш-таблица}: массив $T[0], T[1], \dots, T[m-1]$. Ключ $x \in U$ хранится в ячейке $T[h(x)]$.

Если двум разным ключам $x \ne y$ присвоилось одно и то же значение $h(x) = h(y)$, мы говорим, что произошла \term{коллизия}.
\end{definition}

Иллюстрация работы хеш-функции и появления коллизий приведена на рис.~\ref{fig:p5-1}.

\begin{figure}[H]
\centering
\begin{tikzpicture}[
  >=Stealth, every node/.style={font=\footnotesize},
  uni/.style={draw=prosvBlue, fill=prosvLight, rounded corners, minimum height=0.5cm, minimum width=2.2cm},
  cell/.style={draw=prosvRed!70, fill=prosvCream, minimum height=0.6cm, minimum width=0.7cm},
]
  \node[uni] (k1) at (0, 2)   {ivan@mail};
  \node[uni] (k2) at (0, 1.3) {olga@mail};
  \node[uni] (k3) at (0, 0.6) {petr@yand};
  \node[uni] (k4) at (0, -0.1){anna@bk};
  \node[uni] (k5) at (0, -0.8){vova@mail};

  \node[font=\footnotesize, prosvBlue] at (0, 2.7) {ключи $x$ (вселенная $U$)};

  \foreach \i in {0,1,2,3,4,5,6,7} {
    \node[cell] (c\i) at (5 + 0.75*\i, 0.5) {\i};
  }
  \node[font=\footnotesize, prosvBlue] at (7.5, 1.3) {хеш-таблица $T$ (размер $m=8$)};

  \draw[->, prosvGreen!60!black] (k1.east) -- (c3.west);
  \draw[->, prosvGreen!60!black] (k2.east) -- (c0.west);
  \draw[->, prosvGreen!60!black] (k3.east) -- (c5.west);
  \draw[->, prosvGreen!60!black] (k4.east) -- (c1.west);
  \draw[->, prosvRed!80, very thick] (k5.east) -- (c3.west);

  \node[font=\footnotesize, prosvRed, align=center] at (4, -1.2) {две стрелки в одну\\ ячейку $=$ коллизия};
\end{tikzpicture}
\caption{Хеш-функция «утрамбовывает» большую вселенную ключей в маленькую таблицу. Иногда возникают коллизии.}
\label{fig:p5-1}
\end{figure}

Если коллизий нет (или их мало), мы успешно сжали огромную вселенную ключей в небольшую таблицу. Если же коллизий много, эффективность падает: на одну ячейку приходится много ключей, поиск замедляется.

% ---------------------------------------------------------------------
\subsection{Беда детерминированного хеширования}

Зафиксируем какую-нибудь конкретную хеш-функцию --- например, $h(x) = x \bmod m$ (в предположении, что ключи $x$ представлены как целые числа). На большинстве входов работает прекрасно. Но злоумышленник, зная нашу функцию, может специально подобрать набор ключей $S = \{0, m, 2m, 3m, \dots\}$ --- все они дают одинаковый остаток~$0$. Все они попадут в одну ячейку. Хеш-таблица деградирует до линейного списка, и каждый поиск занимает $O(n)$ вместо $O(1)$.

Этот сценарий не теоретический: реальные сетевые сервисы испытывали атаки именно такого рода --- атаки «алгоритмическим усложнением», когда поток специально подобранных запросов приводил к замедлению сервера в тысячи раз.

\begin{importantbox}[title={Важно: принципиальное наблюдение}]
Для \emph{любой} фиксированной хеш-функции $h$ найдётся «плохой» набор ключей, на котором коллизий очень много. Никакая «гениальная» формула $h$ не спасает.

Спасает другое: \emph{не фиксировать $h$ заранее}, а выбирать её случайно из специального семейства функций.
\end{importantbox}

% ---------------------------------------------------------------------
\subsection{Универсальное семейство хеш-функций}

Идея, восходящая к Картеру и Вегману (1977), состоит в следующем. Договоримся об \emph{целом семействе} хеш-функций $\mathcal{H} = \{h_a\}$. При создании хеш-таблицы выберем из $\mathcal{H}$ \emph{случайную} функцию $h_a$. Эта функция и будет использоваться. Поскольку выбор делается уже после того, как зафиксированы данные (или независимо от них), злоумышленник заранее не знает, какая $h_a$ будет применена, и не может «подстроиться» под неё.

Чтобы это сработало, семейство $\mathcal{H}$ должно обладать специальным свойством.

\begin{definition}{универсальное семейство}{def:p5-2}
Семейство хеш-функций $\mathcal{H} = \{h: U \to \{0, \dots, m-1\}\}$ называется \term{универсальным}, если для любых двух различных ключей $x, y \in U$ ($x \ne y$) вероятность коллизии при случайном выборе $h \in \mathcal{H}$ не превосходит $1/m$:
\[
  \Pr_{h \in \mathcal{H}}[\,h(x) = h(y)\,] \le \frac{1}{m}.
\]
\end{definition}

Поясним: если бы $h$ присваивала каждому ключу совершенно случайный номер ячейки из $m$ возможных, вероятность совпадения двух конкретных ключей в одной ячейке составила бы ровно $1/m$. Универсальное семейство просит лишь того же порядка --- вероятность \emph{не больше} $1/m$.

\begin{theorem}{оценка коллизий}{thm:p5-1}
Если хеш-функция выбрана из универсального семейства, то математическое ожидание числа коллизий заданного ключа $x$ с прочими $n - 1$ ключами таблицы не превосходит $(n-1)/m$. В частности, если $m \ge n$, среднее число коллизий не превосходит $1$.
\end{theorem}

\paragraph*{Доказательство.} Введём индикаторы: $X_y = 1$, если $h(x) = h(y)$, иначе $X_y = 0$. Общее число коллизий ключа $x$ есть $\sum_{y \in S, y \ne x} X_y$. По линейности математического ожидания:
\[
  \mathbb{E}\!\left[\sum X_y\right] = \sum \mathbb{E}[X_y] = \sum \Pr[h(x) = h(y)] \le (n-1) \cdot \frac{1}{m}.
\]
\hfill$\square$

Таким образом, для $m \approx n$ \emph{в среднем} на одну ячейку приходится не более одной--двух коллизий, и время поиска остаётся константным.

\paragraph*{Где встречается случайность.} Подчеркнём: математическое ожидание берётся \emph{по случайному выбору $h$}, а не по случайному выбору ключей! Сами ключи могут быть «худшими из возможных» --- утверждение всё равно справедливо. Это и есть сила рандомизации.

% ---------------------------------------------------------------------
\subsection{Конструкция универсального семейства через теорию чисел}

Теперь самый интересный вопрос: \emph{как же построить} такое семейство $\mathcal{H}$? И вот тут на сцену выходит теория чисел.

\paragraph*{Шаг 1: представление ключа.} Будем считать, что каждый ключ $x \in U$ --- это \emph{вектор} из $k$ небольших чисел: $x = (x_1, x_2, \dots, x_k)$. Например, если ключ --- это IP-адрес вида 192.168.0.1, его естественно представить как четвёрку чисел $(192, 168, 0, 1)$. Если ключ --- это строка, разобьём её на четвёрки или восьмёрки байтов.

Пусть каждое $x_i$ лежит в диапазоне $0 \le x_i < p$, где $p$ --- \term{простое число}, чуть большее размера наших значений.

\paragraph*{Шаг 2: выбор размера таблицы.} Положим $m = p$. (В этом разделе для простоты принимаем такое допущение; в общем случае схема легко обобщается.)

\paragraph*{Шаг 3: семейство $\mathcal{H}$.} Для каждого вектора $a = (a_1, \dots, a_k)$, где $0 \le a_i < p$, определим хеш-функцию:
\[
  h_a(x) = \big( a_1 x_1 + a_2 x_2 + \dots + a_k x_k \big) \bmod p.
\]
Семейство $\mathcal{H} = \{h_a : a \in \{0, \dots, p-1\}^k\}$. Размер семейства --- $p^k$.

\paragraph*{Шаг 4: случайный выбор.} При построении таблицы выбираем $a$ случайно равновероятно из $\{0, \dots, p-1\}^k$ (например, генератором случайных чисел).

\begin{theorem}{Картер--Вегман}{thm:p5-2}
Описанное выше семейство $\mathcal{H}$ \emph{универсально}.
\end{theorem}

\paragraph*{Доказательство.} Возьмём два разных ключа $x = (x_1, \dots, x_k)$ и $y = (y_1, \dots, y_k)$, $x \ne y$. Поскольку $x \ne y$, найдётся хотя бы один индекс $i$, в котором $x_i \ne y_i$; без ограничения общности будем считать $i = 1$.

Коллизия $h_a(x) = h_a(y)$ означает:
\[
  a_1 x_1 + a_2 x_2 + \dots + a_k x_k \equiv a_1 y_1 + a_2 y_2 + \dots + a_k y_k \pmod p,
\]
или, после перестановки,
\[
  a_1 (x_1 - y_1) \equiv -\sum_{i=2}^{k} a_i (x_i - y_i) \pmod p. \qquad (\ast)
\]
Обозначим правую часть как $R$. Заметим, что $R$ зависит от $a_2, \dots, a_k$, но \emph{не зависит} от $a_1$.

\textbf{Здесь и появляется теория чисел.} Поскольку $p$ --- \emph{простое}, а $x_1 - y_1 \ne 0 \pmod p$ (по нашему предположению $x_1 \ne y_1$, и оба они в диапазоне $[0, p)$, значит их разность ненулевая по модулю~$p$), у числа $(x_1 - y_1)$ \emph{существует и единственен} обратный элемент $(x_1 - y_1)^{-1} \bmod p$. Это --- ровно то свойство модулярной арифметики над простым модулем, которое мы доказали в §\ref{sec:numtheory} (расширенный алгоритм Евклида).

Следовательно, уравнение $(\ast)$ имеет ровно одно решение относительно $a_1$ при фиксированных $a_2, \dots, a_k$:
\[
  a_1 = R \cdot (x_1 - y_1)^{-1} \bmod p.
\]
Среди $p$ возможных значений $a_1$ только одно приводит к коллизии. Поэтому условная вероятность коллизии (при фиксированных $a_2, \dots, a_k$) равна $1/p$. Усреднив по $a_2, \dots, a_k$, получаем
\[
  \Pr_a[h_a(x) = h_a(y)] = \frac{1}{p} = \frac{1}{m}.
\]
\hfill$\square$

\begin{importantbox}[title={Важно: почему именно простое $p$}]
Ключевую роль играет то, что $p$ --- \emph{простое}. Будь $p$ составным, у числа $x_1 - y_1$ \emph{могло бы не быть обратного} (если бы оно делилось на общий с $p$ множитель), и тогда уравнение $a_1 \cdot (x_1 - y_1) \equiv R \pmod p$ либо не имело бы решения, либо имело несколько --- и аккуратной оценки $1/p$ не получилось бы.

Так теорема единственности обратного элемента по простому модулю, открытая в античности, оказывается фундаментом современных структур данных, ежесекундно работающих во всех серверах мира.
\end{importantbox}

% ---------------------------------------------------------------------
\subsection{Пример: маленькая хеш-таблица}

Пусть простое $p = 7$ (значит, и $m = 7$), а ключи --- двойки чисел из $\{0, 1, \dots, 6\}$, то есть $k = 2$. Возьмём пять «настоящих» ключей:
\[
  S = \{(1, 2),\ (3, 4),\ (5, 5),\ (2, 6),\ (4, 1)\}.
\]
Случайно выбираем $a = (a_1, a_2)$. Например, пусть случайный генератор выдал $a = (3, 5)$. Тогда $h_a(x) = (3 x_1 + 5 x_2) \bmod 7$:

\begin{center}
\renewcommand{\arraystretch}{1.25}
\begin{tabular}{ccc}
\rowcolor{prosvLight}
ключ $x$ & $3 x_1 + 5 x_2$ & $\bmod 7$ \\
\hline
$(1, 2)$ & $3 + 10 = 13$  & $6$ \\
$(3, 4)$ & $9 + 20 = 29$  & $1$ \\
$(5, 5)$ & $15 + 25 = 40$ & $5$ \\
$(2, 6)$ & $6 + 30 = 36$  & $1$ \\
$(4, 1)$ & $12 + 5 = 17$  & $3$ \\
\end{tabular}
\end{center}

Получаем коллизию: $(3,4)$ и $(2,6)$ попали в ячейку $1$. Если бы случайный генератор выдал другое $a$, например $a = (2, 1)$, картина была бы другой:

\begin{center}
\renewcommand{\arraystretch}{1.25}
\begin{tabular}{ccc}
\rowcolor{prosvLight}
ключ $x$ & $2 x_1 + x_2$ & $\bmod 7$ \\
\hline
$(1, 2)$ & $4$  & $4$ \\
$(3, 4)$ & $10$ & $3$ \\
$(5, 5)$ & $15$ & $1$ \\
$(2, 6)$ & $10$ & $3$ \\
$(4, 1)$ & $9$  & $2$ \\
\end{tabular}
\end{center}

Снова коллизия (правда, в другой ячейке: $(3,4)$ и $(2,6)$ дали $3$). Удивительно: эти два ключа коллидируют при разных $a$. Дело в том, что $(3, 4) - (2, 6) = (1, -2)$, и многие $a$ удовлетворяют $a_1 - 2 a_2 \equiv 0 \pmod 7$. Однако в среднем по всем $a$ вероятность коллизии любых двух конкретных ключей --- ровно $1/7$, как и обещает теорема.

% ---------------------------------------------------------------------
\subsection{Применения хеширования}

Хеширование --- одна из самых востребованных идей в современной информатике. \emph{Универсальное} хеширование особенно важно для хеш-таблиц и быстрого поиска, а \emph{криптографические} хеш-функции --- для проверки целостности, цифровой подписи и блокчейна. Вот несколько мест, где работают эти идеи:

\begin{itemize}
  \item \term{Поиск дубликатов и быстрая выборка.} Все базы данных, все таблицы соответствий в операционных системах, все словари в языках программирования (Python: \texttt{dict}, JavaScript: \texttt{Map}, Java: \texttt{HashMap}) построены на хешировании.
  \item \term{Проверка целостности файлов.} Контрольные суммы при скачивании файла (CRC, MD5, SHA) --- по сути, хеш-функции. Если хеш скачанного файла не совпадает с заявленным, файл повреждён.
  \item \term{Цифровая подпись.} Как мы видели в §\ref{subsec:eds}, RSA-подпись на самом деле ставится не на весь документ, а на его хеш.
  \item \term{Антивирусы и системы обнаружения дубликатов.} База вирусных сигнатур --- это хеши известных вредоносных файлов; чтобы проверить, не заражён ли файл, антивирус считает его хеш и ищет в базе.
  \item \term{Block\-chain.} В криптовалютах хеш каждого блока включается в следующий блок, образуя цепочку, в которой невозможно незаметно подправить «прошлое».
\end{itemize}

\begin{interestbox}[title={Это интересно: криптографические хеш-функции}]
В этом параграфе мы говорили о «простых» хеш-функциях, минимизирующих коллизии. В криптографии используется родственное, но более сильное понятие --- \term{криптографическая хеш-функция}. От неё требуется не только малое число коллизий, но и \emph{вычислительная невозможность найти} два разных входа с одинаковым хешем (\engterm{collision resistance}), а также найти прообраз заданного хеша (\engterm{preimage resistance}). Примеры таких функций --- SHA-256, BLAKE2, в России --- «Стрибог» (ГОСТ Р 34.11--2012). Внутри они устроены сложнее (и не на основе модулярной арифметики), но используют тот же общий принцип: смешать входные биты так, чтобы малейшее изменение во входе приводило к радикальному изменению в выходе.
\end{interestbox}

\begin{tasksbox}
\begin{enumerate}
  \item Для простого $p = 5$, $k = 1$ и $a = 3$ найдите $h_a(x) = a x \bmod p$ для всех $x \in \{0, 1, 2, 3, 4\}$.
  \item Пусть ключи --- пары $(x_1, x_2)$, $p = 11$, $a = (4, 7)$. Найдите коллидирует ли пара $(1, 3)$ с парой $(8, 1)$ при таком $a$.
  \item Объясните, что произойдёт с оценкой $1/m$, если в семействе Картера--Вегмана взять $p$ составным (например, $p = 6$).
  \item Почему в формуле хеш-функции $h_a(x) = (a_1 x_1 + \dots + a_k x_k) \bmod p$ нет «свободного члена» $a_0$, как в линейной функции $a x + b$? Что изменится, если его добавить?
  \item* Покажите, что если выбирать $a$ всего лишь из множества $\{1\}$ (одно значение), то семейство не универсально.
  \item* Покажите, что вероятность того, что у заданного ключа \emph{не будет} коллизий ни с каким из $n-1$ других ключей в таблице размера $m \ge n$, не меньше $1/2$ (используйте неравенство Маркова).
\end{enumerate}
\end{tasksbox}
